\documentclass[11pt]{article}
\usepackage[margin=2.2cm]{geometry}
\usepackage{amsmath,amssymb,amsthm,mathtools}
\usepackage{algorithm}
\usepackage{algpseudocode}
\usepackage{booktabs}
\usepackage[hidelinks]{hyperref}
\usepackage{tikz}
\usetikzlibrary{shapes.geometric,arrows.meta,positioning,fit,backgrounds}

\newtheorem{theorem}{Theorem}
\newtheorem{definition}{Definition}
\newtheorem{remark}{Remark}

\newcommand{\Ks}{\mathcal{K}^{\star}}
\newcommand{\thl}{\underline{\theta}}
\newcommand{\thu}{\overline{\theta}}

\title{HICA-S, Track T4: Three Pruning Layers of Algorithm A\\
\large Label dominance, per-bundle Pareto filtering, and the FD-completion label rule}
\author{Tran An Khoa}
\date{September 2026}

\begin{document}
\maketitle

\begin{abstract}
Algorithm A (route generation for one driver) enumerates every feasible pickup/delivery
sequence and reduces it to a route pool using three successive, independently safe pruning
layers, followed by a fourth, separate post-processing filter ($\Ks$) that operates on the
finished pool. This note gives each layer its own precise statement, an algorithm block, a
worked numeric example, and a flowchart showing how they compose into the full Algorithm~A
$\to$ $\Ks$ pipeline. Layer 3 (the FD-completion label rule) and $\Ks$ share the same
underlying inequality --- ``a sub-bundle route plus the FD price for what is dropped can beat
the full route'' --- but apply it at different times: Layer 3 applies a \emph{lower bound} to
an \emph{unfinished} label so that hopeless subtrees are never expanded; $\Ks$ applies the
\emph{exact} inequality to \emph{finished} routes so that hopeless entries are removed from
the pool before Algorithms B/C ever see them. This is why enabling Layer 3 never changes the
pool that $\Ks$ eventually returns (empirically confirmed: 0 violations across 5{,}654{,}551
audited completions on 5 real RQ1 instances, $B\in\{3,4\}$).
\end{abstract}

\tableofcontents

%=====================================================================
\section{Where each layer sits in the pipeline}
%=====================================================================

Layers 1--3 all execute \emph{inside} the depth-first label-extension recursion of
Algorithm~A, for a single driver $i$. $\Ks$ executes \emph{after} Algorithm~A has returned a
finished pool $R_i$, before that pool is handed to Algorithm~B (the winner-determination
MILP) and Algorithm~C (the VCG removal solves). Figure~\ref{fig:flow} gives the full picture;
Table~\ref{tab:overview} gives a one-line comparison.

\begin{table}[h]
\centering
\begin{tabular}{@{}lllll@{}}
\toprule
Layer & Compares & Scope & Uses FD price? & Needs completion? \\
\midrule
1 (dominance)      & label vs.\ label        & same $(v,IV,C)$        & no  & no  \\
2 (Pareto/bundle)   & route vs.\ route        & same bundle $S$        & no  & yes \\
3 (FD-completion)   & label vs.\ FD           & any $T\subseteq P(L)$  & yes & no (lower bound) \\
$\Ks$ (post-filter) & route vs.\ FD           & any $T\subseteq S$     & yes & yes (exact) \\
\bottomrule
\end{tabular}
\caption{The four filters, in the order they can fire. Layers 1--3 run during route
generation (Algorithm~A); $\Ks$ runs once, after Algorithm~A returns.}
\label{tab:overview}
\end{table}

%=====================================================================
\section{Layer 1 --- Label dominance (Dumas, Desrosiers \& Soumis, 1991)}
%=====================================================================

\subsection{What question it answers}
While a route is still being built, many partial paths (\emph{labels}) may reach the same
``state'' by different routes so far. Layer 1 asks: among labels that share a state, which
ones can already be discarded because another label is at least as good on every axis that
matters for the future?

\begin{definition}[Label state and dominance]\label{def:label-dom}
A label is a tuple $L=(v,IV,C,t,K,W)$: current stop $v$, orders on board $IV$, orders fully
served $C$, current time $t$, and accumulated raw cost $(K,W)$. Two labels $L_a,L_b$ are
\emph{comparable} iff $(v_a,IV_a,C_a)=(v_b,IV_b,C_b)$. Label $L_a$ \emph{dominates} $L_b$
(written $L_a\preceq L_b$) iff
\[
t_a\le t_b,\qquad K_a\le K_b,\qquad W_a\le W_b,
\]
with at least one inequality strict.
\end{definition}

\begin{remark}[Why this is always safe]
If $L_a\preceq L_b$, then any completion $\sigma$ feasible from $L_b$ is also feasible from
$L_a$ (same on-board set, same served set, same or earlier position in time cannot make a
downstream deadline harder to meet), and the resulting route from $L_a\cdot\sigma$ has cost
component-wise no larger than from $L_b\cdot\sigma$. Hence $L_b$ can never lead to a route that
strictly beats every route reachable from $L_a$; deleting $L_b$ and its whole subtree loses
nothing.
\end{remark}

\begin{algorithm}[H]
\caption{Layer 1 --- filter dominated labels at one DP round}
\label{alg:layer1}
\begin{algorithmic}[1]
\Function{FilterDominated}{labels}
  \State $n \gets |\text{labels}|$; $\text{dominated} \gets [\text{False}]\times n$
  \For{$i \gets 1$ \textbf{to} $n$}
    \If{dominated$[i]$} \textbf{continue} \EndIf
    \For{$j \gets 1$ \textbf{to} $n$, $j\ne i$}
      \If{dominated$[j]$} \textbf{continue} \EndIf
      \If{labels$[j] \preceq$ labels$[i]$} \State dominated$[i]\gets$ True \State \textbf{break} \EndIf
      \If{labels$[i] \preceq$ labels$[j]$} \State dominated$[j]\gets$ True \EndIf
    \EndFor
  \EndFor
  \State \Return $\{\text{labels}[i] : \neg\text{dominated}[i]\}$
\EndFunction
\end{algorithmic}
\end{algorithm}

\paragraph{Worked example.} Two labels reach state $(v{=}n_8,\ IV{=}\{o_1,o_3\},\
C{=}\{o_5\})$: label $A=(t{=}42,\ K{=}5.1,\ W{=}0.6)$ and label $B=(t{=}47,\ K{=}5.1,\
W{=}0.8)$. Since $t_A\le t_B$, $K_A\le K_B$ (equal), $W_A\le W_B$ with the last inequality
strict, $A\preceq B$: $B$ is deleted, $A$ survives. Any completion later tried from $A$ starts
no later and no more expensive than the same completion from $B$ would have been.

\paragraph{Where it lives in the code.} \texttt{t6\_dp.py::\_filter\_dominated\_labels},
re-applied after every DP round (mirrored for audit purposes in
\texttt{hica\_core.py::\_filter\_dominated\_labels}). It is present in \emph{every} run of
Algorithm~A, whether or not Layer 3 is enabled --- it is not a variable in any of the
before/after comparisons reported elsewhere in this track.

%=====================================================================
\section{Layer 2 --- Per-bundle Pareto filtering}
%=====================================================================

\subsection{What question it answers}
Once Algorithm~A has finished exploring, several \emph{complete} routes may serve exactly the
same bundle of orders $S$ (same set, different stop order, different amount of detour). Layer
2 asks: among the finished routes that all serve exactly $S$, which are dominated purely on
$(K,W)$?

\begin{definition}[Per-bundle Pareto dominance]\label{def:layer2}
Fix a bundle $S$ and let $R_S=\{r_1,\dots,r_m\}$ be every complete route generated by
Algorithm~A with bundle exactly $S$. Route $r_a$ \emph{Pareto-dominates} $r_b$ (both in $R_S$)
iff $K_{r_a}\le K_{r_b}$ and $W_{r_a}\le W_{r_b}$, with at least one strict. The
\emph{per-bundle Pareto front} kept for $S$ is
\[
\mathrm{PF}(S) \;=\; \{\, r\in R_S : \nexists\, r'\in R_S,\ r' \text{ Pareto-dominates } r \,\}.
\]
\end{definition}

\begin{remark}[Why the comparison never crosses bundles]
Every $r\in R_S$ serves the \emph{identical} order set $S$; hence for any bid $b$,
$c_r(b)=K_r+bW_r$ and $c_{r'}(b)=K_{r'}+bW_{r'}$ are directly comparable at every $b$ without
any correction term. If $r'$ Pareto-dominates $r$, then $c_{r'}(b)\le c_r(b)$ for \emph{every}
$b\ge0$, so $r$ can never be strictly cheaper than $r'$ and is safe to discard. This is
different in kind from comparing routes with \emph{different} bundles $S\ne S'$: there the
``missing'' orders $S\setminus S'$ must go somewhere (FD, or another driver), and a raw
$(K,W)$ comparison that ignores this is unsound (Remark~\ref{rem:hull-rejected}).
\end{remark}

\begin{algorithm}[H]
\caption{Layer 2 --- Pareto front of one bundle}
\label{alg:layer2}
\begin{algorithmic}[1]
\Function{ParetoFront}{$\text{kw\_list} = [(K_1,W_1),\dots,(K_m,W_m)]$}
  \State sort kw\_list by $(K,W)$
  \State keep $\gets [\,]$; bestW $\gets +\infty$
  \For{$(K,W)$ in sorted kw\_list}
    \If{$W <$ bestW} \State append $(K,W)$ to keep \State bestW $\gets W$ \EndIf
  \EndFor
  \State \Return keep
\EndFunction
\end{algorithmic}
\end{algorithm}

\paragraph{Worked example.} Bundle $S=\{o_2,o_7,o_9\}$ has three complete routes found by
Algorithm~A: $A=(K{=}10,\,W{=}1.2)$, $B=(K{=}12,\,W{=}0.9)$, $C=(K{=}15,\,W{=}1.5)$. $C$ is
dominated by $A$ on both axes ($10\le15$, $1.2\le1.5$) and discarded. $A$ and $B$ each win on
one axis ($A$ cheaper in $K$, $B$ cheaper in $W$), so both survive: which one is optimal
depends on the bid $b$ that will be revealed later.

\begin{remark}[What Layer 2 is \emph{not}]\label{rem:hull-rejected}
An earlier, rejected proposal (\texttt{Report\_ActivationRate.md}, \S5) tried to build a
\emph{single} Pareto/convex hull over the \emph{entire} pool of one driver, mixing routes from
every bundle. That rule is unsound: a route of a larger bundle can still be the unique optimum
at some bid even though it is not on the driver-wide hull, because serving the extra orders
\emph{together} exploits economies of scale that a per-bundle hull cannot see, and because the
alternative for the missing orders (FD, or another driver already busy elsewhere) is not
accounted for by a raw geometric hull. That proposal was audited and discarded (4 of 6
genuinely activated routes fell outside its candidate set on the very first smoke test); Layer
2, restricted to \emph{one bundle at a time}, does not have this flaw and is unconditionally
safe (Definition~\ref{def:layer2}).
\end{remark}

\paragraph{Where it lives in the code.} \texttt{dp\_labeling.py::\_pareto\_front}, called once
per finished bundle (\texttt{for Cset, labs in complete\_by\_C.items(): front =
\_pareto\_front(kw)}) --- never across bundles.

%=====================================================================
\section{Layer 3 --- The FD-completion label rule}
%=====================================================================

\subsection{What question it answers}
Layers 1--2 only ever compare entities \emph{internal} to the driver's own search (label vs.\
label, or finished route vs.\ finished route of the same bundle). Layer 3 is the first layer
to compare against an \emph{external} price: the fixed-fleet (FD) tariff $q(\cdot)$. It asks,
mid-construction: \emph{having already served a set $T$ of orders in the current label, is it
already certain --- for every way the route might be finished, and every admissible bid --
that keeping $T$ on board costs more than simply handing $T$ to FD?}

\begin{definition}[Shortcut and its lower-bound savings, informal]\label{def:layer3}
For a label $L$ with served set $T\subseteq P(L)$, let $L\setminus T$ be $L$ with $T$'s stops
removed (the \emph{shortcut}). Define
\[
\Delta D(L,T) \;=\; \text{distance $L$ spends on $T$ (junction-corrected)}, \qquad
\delta(L,T) \;=\; \text{time $L$ spends on $T$ (junction-corrected)},
\]
\[
A(L,T) \;=\; \text{worst-case time the shortcut's later \emph{waiting} could absorb},
\qquad
\beta(L,T) \;=\; \kappa\,\Delta D \;+\; \thl\,\frac{\max\{0,\,\delta-A\}}{60}.
\]
\emph{Rule:} discard $L$ (and its entire subtree) if $\beta(L,T)\ge q(T)$ for some
$\emptyset\ne T\subseteq P(L)$.
\end{definition}

\begin{remark}[Why $\beta$ must be a \emph{lower} bound, and why $A$ is unavoidable]
$L$ is unfinished, so the true saving from dropping $T$ depends on how the route is completed
--- unknown at this point. $\beta(L,T)$ is proved (Theorem~1 of \texttt{T4\_Label\_Rule.tex})
to be a valid lower bound on that saving, for \emph{every} feasible completion and every bid
$b\ge\thl$. Naively transplanting rollback pruning from the ESPPRC literature (Lozano, Duque \&
Medaglia, 2016) corresponds to setting $A\equiv0$ and is \emph{unsafe} here, because unlike an
additive-cost shortest-path problem, this route also allows \emph{waiting}: a time saving
credited now can be silently absorbed by waiting at a later, not-yet-open pickup, so the
uncorrected rule can discard a route that a full simulation would have kept (a concrete
counterexample with numbers is in \texttt{T4\_Label\_Rule.tex}, Example~1). $A(L,T)$ subtracts
exactly the part of the credited saving that later waiting could eat.
\end{remark}

\begin{algorithm}[H]
\caption{Layer 3 --- look-back FD-completion check (fires right after serving $T$)}
\label{alg:layer3}
\begin{algorithmic}[1]
\Function{TryFire}{$L$, served set $T$, FD prices $q$, $\thl$, $\kappa$}
  \State $q_T \gets \sum_{j\in T} q_j$
  \State $(D',v',r') \gets$ schedule of the shortcut $L\setminus T$ \Comment{junction-corrected}
  \State $\Delta D \gets D_L - D' - d(v',v_L)$
  \State $\delta \gets t_L - r' - \tau(v',v_L)$
  \State $\mathcal{F} \gets$ future reachable stops from $(v_L,t_L)$ \Comment{§\ref{def:layer3} remark on deliveries}
  \State $A \gets \max\bigl(0,\ \max_{w\in\mathcal{F}}(e_w - r' - \tau(v',w))\bigr)$
  \State $\beta \gets \kappa\,\Delta D + \thl\cdot\max(0,\delta-A)/60$
  \If{$\beta \ge q_T$}
    \State \textbf{prune} $L$ and its entire subtree \Comment{safe by Theorem 1, \texttt{T4\_Label\_Rule.tex}}
  \EndIf
\EndFunction
\end{algorithmic}
\end{algorithm}

\paragraph{Worked example (numbers from Example~1 of \texttt{T4\_Label\_Rule.tex}).} Label
$L=(p_{o_2},d_{o_2})$, $T=\{o_2\}$, $q_{o_2}=17.13$. Ignoring absorption ($A\equiv0$, the
unsafe transplant) credits a $42.9$-minute saving, giving $\beta=17.72\ge17.13$: the label
would be wrongly discarded. With $A$ correctly computed, the true saving at $\thl$ is only
$7.24<17.13$: the label is (correctly) kept, because the true completion
$(p_{o_2},d_{o_2},p_{o_0},d_{o_0})$ needs to wait for $o_0$'s late opening anyway, absorbing
most of the naive saving.

\paragraph{Where it lives in the code.} \texttt{dp\_rules.py} (full enumerator, used for
correctness auditing) and \texttt{dp\_fast.py} (incremental $O(1)$-per-extension
implementation, used for timing) --- both under the \texttt{'LB2'} rule name. It requires no
completed route: it fires purely from information already accumulated in the label plus a
$O(1)$-updated shortcut state per picked order.

%=====================================================================
\section{$\Ks$ --- the exact post-processing filter (not a fourth DP layer)}
%=====================================================================

\begin{definition}[Local pruning frontier $\Ks_i$]\label{def:kstar}
For driver $i$'s \emph{finished} pool $R_i$ (after Layers 1--2), and for a route $r\in R_i$
with bundle $S_r$, define the envelope of everything that could beat $r$ at bid $b$:
\[
E_r(b) \;=\; \min\Bigl\{\, q(S_r),\ \min_{\substack{r'\in R_i,\ r'\ne r\\ S_{r'}\subsetneq S_r}}
\bigl(c_{r'}(b) + q(S_r\setminus S_{r'})\bigr) \Bigr\}, \qquad
\mu_r(b) \;=\; E_r(b) - c_r(b).
\]
Then $\Ks_i = \{\, r\in R_i : \max_{b\in[\thl,\thu]}\mu_r(b) > 0 \,\}$: routes that can win
against \emph{every} sub-bundle alternative (including sending the rest to FD) for at least
one admissible bid.
\end{definition}

\begin{remark}[Same inequality as Layer 3, different time and different scope]
Both Layer 3 and $\Ks$ test the same underlying fact --- ``a piece of this route, plus FD for
the rest, can beat the whole route.'' The differences: (i) $\Ks$ compares against \emph{every}
sub-bundle $T\subseteq S_r$ of the \emph{finished} route, using its \emph{exact} $(K,W)$;
Layer~3 only compares against the set $T$ \emph{already served so far} in an \emph{unfinished}
label, using a provable \emph{lower bound} $\beta\le$ true saving. (ii) $\Ks$ runs once, after
Algorithm~A returns, directly shrinking the pool handed to Algorithms B/C; Layer~3 runs
\emph{during} generation, only saving the work of enumerating routes that $\Ks$ would have
discarded anyway. Theorem~2 of \texttt{T4\_Label\_Rule.tex} proves
$\Ks_i(R_i)=\Ks_i(R'_i)$ where $R'_i$ is the pool generated \emph{with} Layer~3 --- i.e.\
turning Layer~3 on never changes what $\Ks$ eventually returns. This was also confirmed
empirically (Table~\ref{tab:audit}).
\end{remark}

\begin{table}[h]
\centering
\begin{tabular}{@{}lrrr@{}}
\toprule
$B$ & completions audited & bound violations & $\Ks$ changed on any driver? \\
\midrule
3 & 114{,}627 & 0 & no \\
4 & 5{,}539{,}924 & 0 & no \\
\bottomrule
\end{tabular}
\caption{Empirical confirmation on the 5 real RQ1 instances used throughout this track
(\texttt{audit\_rq1\_bounds.py}).}
\label{tab:audit}
\end{table}

%=====================================================================
\section{Flowchart}
%=====================================================================

\begin{figure}[h]
\centering
\resizebox{!}{0.82\textheight}{%
\begin{tikzpicture}[
  node distance=6mm and 10mm,
  every node/.style={font=\footnotesize},
  block/.style={rectangle, rounded corners, draw, fill=blue!6, align=center,
                minimum width=30mm, minimum height=7mm, inner sep=1.5mm},
  decision/.style={diamond, draw, fill=orange!10, align=center, aspect=2.6,
                   minimum width=22mm, inner sep=0.5mm},
  filter/.style={rectangle, draw, fill=green!8, align=center,
                 minimum width=34mm, minimum height=7mm, inner sep=1.5mm},
  final/.style={rectangle, rounded corners, draw, fill=red!8, align=center,
                minimum width=32mm, minimum height=7mm, inner sep=1.5mm},
  lbl/.style={font=\scriptsize},
  arr/.style={-{Stealth[length=1.8mm]}, thick},
  every fit/.style={draw, dashed, rounded corners, inner sep=3mm}
]

% ---- main vertical spine, single column, no left/right zigzag ----
\node[block] (start) {start label\\$(v_0,\emptyset,\emptyset,t_0,0,0)$};
\node[decision, below=of start] (extend) {extend to\\next stop};
\node[filter, below=of extend] (l3) {\textbf{Layer 3}: $\beta(L,T)\ge q(T)$\\for some $T$?};
\node[decision, below=of l3] (fire3) {fires?};
\node[decision, below=of fire3] (complete) {label\\complete?};
\node[filter, below=of complete] (l1) {\textbf{Layer 1}: dominated by\\same-state label?};
\node[decision, below=of l1] (dom1) {dominated?};
\node[block, below=of dom1] (route) {emit route $(S,K,W)$\\into $R_S$};
\node[filter, below=of route] (l2) {\textbf{Layer 2}: per-bundle\\Pareto filter on $R_S$};
\node[block, below=of l2] (pool) {finished pool $R_i$\\(after Layers 1--2)};
\node[filter, below=of pool] (kstar) {$\Ks$: exact post-filter\\(Def.~\ref{def:kstar})};
\node[final, below=of kstar] (out) {$\Ks_i\subseteq R_i$ handed\\to Algorithm B/C};

% ---- side nodes for the two "prune/discard" outcomes, kept OFF the spine ----
\node[block, right=16mm of fire3] (prune3) {prune label\\\& subtree};
\node[block, right=16mm of dom1] (prune1) {discard\\label};

% ---- spine arrows, strictly top-to-bottom ----
\draw[arr] (start) -- (extend);
\draw[arr] (extend) -- (l3);
\draw[arr] (l3) -- (fire3);
\draw[arr] (fire3) -- node[lbl,left]{no} (complete);
\draw[arr] (complete) -- node[lbl,left]{no: extend more} (l1);
\draw[arr] (l1) -- (dom1);
\draw[arr] (dom1) -- node[lbl,left]{no} (route);
\draw[arr] (route) -- (l2);
\draw[arr] (l2) -- (pool);
\draw[arr] (pool) -- (kstar);
\draw[arr] (kstar) -- (out);

% ---- side branches: go right to the prune/discard box, then straight back
%      into the spine at the node ABOVE where the loop restarts (no crossing) ----
\draw[arr] (fire3) -- node[lbl,above]{yes} (prune3);
\draw[arr] (prune3) |- ($(extend.east)+(0,-1mm)$);
\draw[arr] (dom1) -- node[lbl,above]{yes} (prune1);
\draw[arr] (prune1) |- ($(l1.east)+(0,4mm)$);

% ---- "complete=yes" skips Layer 1, goes straight down into "route";
%      routed via an explicit waypoint further right than prune3/prune1,
%      so the two 90-degree turns are unambiguous and it cannot cross them ----
\coordinate (bypass) at ($(complete.east)+(32mm,0)$);
\draw[arr] (complete.east) -- node[lbl,above]{yes} (bypass) -- (bypass |- route.east) -- (route.east);

\begin{scope}[on background layer]
\node[fit=(l3)(fire3)(complete)(l1)(dom1)(prune3)(prune1),
      label={[lbl]above:Algorithm A -- inside the label-extension recursion}] {};
\end{scope}

\end{tikzpicture}%
}
\caption{Full pipeline, read top to bottom. Layers 1 and 3 sit inside the recursive
label-extension loop of Algorithm~A (dashed box): Layer~3 is checked first, right after a
set $T$ has just been served; if the label survives and later becomes complete, Layer~1
dominance is checked against other labels sharing the same state before the route is
emitted. Layer~2 and $\Ks$ run once the whole pool is finished. Both ``prune/discard''
outcomes are drawn off the main spine and loop back into the recursion (no crossing lines).}
\label{fig:flow}
\end{figure}

%=====================================================================
\section{Summary}
%=====================================================================

\begin{itemize}
\item \textbf{Layer 1} (label dominance, Dumas et al.\ 1991) --- always on, compares
      \emph{unfinished} labels sharing the same $(v,IV,C)$, no FD price involved.
\item \textbf{Layer 2} (per-bundle Pareto filter) --- always on, compares \emph{finished}
      routes that serve the \emph{identical} bundle, no FD price involved. Not to be confused
      with the rejected driver-wide hull (Remark~\ref{rem:hull-rejected}).
\item \textbf{Layer 3} (FD-completion label rule, this track's new contribution) --- optional,
      compares an \emph{unfinished} label's served set $T$ against $q(T)$ using a proven lower
      bound $\beta(L,T)$; corrects for waiting absorption ($A$), which a direct transplant of
      rollback pruning omits and which is provably unsafe to omit here.
\item \textbf{$\Ks$} (local pruning frontier) --- a separate, exact post-processing filter
      applied once to the finished pool $R_i$; uses the same inequality as Layer~3 but with
      exact costs and full sub-bundle coverage. Enabling Layer~3 provably (Theorem~2) and
      empirically (0/5{,}654{,}551 violations) never changes $\Ks_i$.
\end{itemize}

\end{document}
